% Options for packages loaded elsewhere
\PassOptionsToPackage{unicode}{hyperref}
\PassOptionsToPackage{hyphens}{url}
\PassOptionsToPackage{dvipsnames,svgnames,x11names}{xcolor}
%
\documentclass[
  letterpaper,
  DIV=11,
  numbers=noendperiod]{scrartcl}

\usepackage{amsmath,amssymb}
\usepackage{iftex}
\ifPDFTeX
  \usepackage[T1]{fontenc}
  \usepackage[utf8]{inputenc}
  \usepackage{textcomp} % provide euro and other symbols
\else % if luatex or xetex
  \usepackage{unicode-math}
  \defaultfontfeatures{Scale=MatchLowercase}
  \defaultfontfeatures[\rmfamily]{Ligatures=TeX,Scale=1}
\fi
\usepackage{lmodern}
\ifPDFTeX\else  
    % xetex/luatex font selection
\fi
% Use upquote if available, for straight quotes in verbatim environments
\IfFileExists{upquote.sty}{\usepackage{upquote}}{}
\IfFileExists{microtype.sty}{% use microtype if available
  \usepackage[]{microtype}
  \UseMicrotypeSet[protrusion]{basicmath} % disable protrusion for tt fonts
}{}
\makeatletter
\@ifundefined{KOMAClassName}{% if non-KOMA class
  \IfFileExists{parskip.sty}{%
    \usepackage{parskip}
  }{% else
    \setlength{\parindent}{0pt}
    \setlength{\parskip}{6pt plus 2pt minus 1pt}}
}{% if KOMA class
  \KOMAoptions{parskip=half}}
\makeatother
\usepackage{xcolor}
\setlength{\emergencystretch}{3em} % prevent overfull lines
\setcounter{secnumdepth}{-\maxdimen} % remove section numbering
% Make \paragraph and \subparagraph free-standing
\makeatletter
\ifx\paragraph\undefined\else
  \let\oldparagraph\paragraph
  \renewcommand{\paragraph}{
    \@ifstar
      \xxxParagraphStar
      \xxxParagraphNoStar
  }
  \newcommand{\xxxParagraphStar}[1]{\oldparagraph*{#1}\mbox{}}
  \newcommand{\xxxParagraphNoStar}[1]{\oldparagraph{#1}\mbox{}}
\fi
\ifx\subparagraph\undefined\else
  \let\oldsubparagraph\subparagraph
  \renewcommand{\subparagraph}{
    \@ifstar
      \xxxSubParagraphStar
      \xxxSubParagraphNoStar
  }
  \newcommand{\xxxSubParagraphStar}[1]{\oldsubparagraph*{#1}\mbox{}}
  \newcommand{\xxxSubParagraphNoStar}[1]{\oldsubparagraph{#1}\mbox{}}
\fi
\makeatother


\providecommand{\tightlist}{%
  \setlength{\itemsep}{0pt}\setlength{\parskip}{0pt}}\usepackage{longtable,booktabs,array}
\usepackage{calc} % for calculating minipage widths
% Correct order of tables after \paragraph or \subparagraph
\usepackage{etoolbox}
\makeatletter
\patchcmd\longtable{\par}{\if@noskipsec\mbox{}\fi\par}{}{}
\makeatother
% Allow footnotes in longtable head/foot
\IfFileExists{footnotehyper.sty}{\usepackage{footnotehyper}}{\usepackage{footnote}}
\makesavenoteenv{longtable}
\usepackage{graphicx}
\makeatletter
\newsavebox\pandoc@box
\newcommand*\pandocbounded[1]{% scales image to fit in text height/width
  \sbox\pandoc@box{#1}%
  \Gscale@div\@tempa{\textheight}{\dimexpr\ht\pandoc@box+\dp\pandoc@box\relax}%
  \Gscale@div\@tempb{\linewidth}{\wd\pandoc@box}%
  \ifdim\@tempb\p@<\@tempa\p@\let\@tempa\@tempb\fi% select the smaller of both
  \ifdim\@tempa\p@<\p@\scalebox{\@tempa}{\usebox\pandoc@box}%
  \else\usebox{\pandoc@box}%
  \fi%
}
% Set default figure placement to htbp
\def\fps@figure{htbp}
\makeatother

\KOMAoption{captions}{tableheading}
\makeatletter
\@ifpackageloaded{caption}{}{\usepackage{caption}}
\AtBeginDocument{%
\ifdefined\contentsname
  \renewcommand*\contentsname{Table des matières}
\else
  \newcommand\contentsname{Table des matières}
\fi
\ifdefined\listfigurename
  \renewcommand*\listfigurename{Liste des Figures}
\else
  \newcommand\listfigurename{Liste des Figures}
\fi
\ifdefined\listtablename
  \renewcommand*\listtablename{Liste des Tables}
\else
  \newcommand\listtablename{Liste des Tables}
\fi
\ifdefined\figurename
  \renewcommand*\figurename{Figure}
\else
  \newcommand\figurename{Figure}
\fi
\ifdefined\tablename
  \renewcommand*\tablename{Table}
\else
  \newcommand\tablename{Table}
\fi
}
\@ifpackageloaded{float}{}{\usepackage{float}}
\floatstyle{ruled}
\@ifundefined{c@chapter}{\newfloat{codelisting}{h}{lop}}{\newfloat{codelisting}{h}{lop}[chapter]}
\floatname{codelisting}{Listing}
\newcommand*\listoflistings{\listof{codelisting}{Liste des Listings}}
\makeatother
\makeatletter
\makeatother
\makeatletter
\@ifpackageloaded{caption}{}{\usepackage{caption}}
\@ifpackageloaded{subcaption}{}{\usepackage{subcaption}}
\makeatother

\ifLuaTeX
\usepackage[bidi=basic]{babel}
\else
\usepackage[bidi=default]{babel}
\fi
\babelprovide[main,import]{french}
% get rid of language-specific shorthands (see #6817):
\let\LanguageShortHands\languageshorthands
\def\languageshorthands#1{}
\usepackage{bookmark}

\IfFileExists{xurl.sty}{\usepackage{xurl}}{} % add URL line breaks if available
\urlstyle{same} % disable monospaced font for URLs
\hypersetup{
  pdftitle={Calcul de l'Autocovariance et des Autocorrélations},
  pdfauthor={Annexes},
  pdflang={fr},
  colorlinks=true,
  linkcolor={blue},
  filecolor={Maroon},
  citecolor={Blue},
  urlcolor={Blue},
  pdfcreator={LaTeX via pandoc}}


\title{Calcul de l'Autocovariance et des Autocorrélations}
\author{Annexes}
\date{Invalid Date}

\begin{document}
\maketitle


\section{Calcul de l'Autocovariance et des
Autocorrélations}\label{calcul-de-lautocovariance-et-des-autocorruxe9lations}

Cette annexe couvre la dérivation de la fonction d'autocorrélation (ACF)
pour les modèles MA(1), MA(Q), AR(1), AR(2), AR(3) et ARMA(1,1). Tout au
long de cette annexe, \(\{ \varepsilon_t \}\) est supposé être un
processus de bruit blanc, et les paramètres des processus sont toujours
supposés être cohérents avec la stationnarité de covariance. Tous les
modèles sont supposés avoir une moyenne de 0, une hypothèse faite sans
perte de généralité puisque les autocovariances sont définies en
utilisant des séries temporelles décentrées,

\[
\gamma_s = E[(Y_t - \mu)(Y_{t-s} - \mu)]
\]

où \(\mu = E[Y_t]\). Rappelons que l'autocorrélation est simplement le
rapport de l'autocovariance d'ordre \(s\) à la variance,

\[
\rho_s = \frac{\gamma_s}{\gamma_0}.
\]

Cette annexe présente deux méthodes pour dériver les autocorrélations
des processus ARMA : la substitution arrière et les équations de
Yule-Walker, un ensemble de \(k\) équations avec \(k\) inconnues où
\(\gamma_0, \gamma_1, \ldots, \gamma_{k-1}\) sont la solution.

\subsection{A.1 Équations de
Yule-Walker}\label{a.1-uxe9quations-de-yule-walker}

Les équations de Yule-Walker sont un système linéaire de
\(\max(P, Q) + 1\) équations (dans un ARMA(P, Q)) où la solution du
système est la variance à long terme et les premières \(k-1\)
autocovariances. Les équations de Yule-Walker sont formées en égalisant
la définition d'une autocovariance avec une expansion produite en
substituant la valeur contemporaine de \(Y_t\). Par exemple, supposons
que \(Y_t\) suit un processus AR(2),

\[
Y_t = \phi_1 Y_{t-1} + \phi_2 Y_{t-2} + \varepsilon_t
\]

La variance doit satisfaire

\[
\begin{aligned}
E[Y_t Y_t] &= E[Y_t (\phi_1 Y_{t-1} + \phi_2 Y_{t-2} + \varepsilon_t)] \\
E[Y_t^2] &= E[\phi_1 Y_t Y_{t-1} + \phi_2 Y_t Y_{t-2} + Y_t \varepsilon_t] \\
V[Y_t] &= \phi_1 E[Y_t Y_{t-1}] + \phi_2 E[Y_t Y_{t-2}] + E[Y_t \varepsilon_t].
\end{aligned}
\]

Dans l'équation finale ci-dessus, les termes de la forme
\(E[Y_t Y_{t-s}]\) sont remplacés par leurs valeurs de population,
\(\gamma_s\), et \(E[Y_t \varepsilon_t]\) est remplacé par sa valeur de
population, \(\sigma^2\).

\[
\begin{aligned}
V[Y_t Y_t] &= \phi_1 E[Y_t Y_{t-1}] + \phi_2 E[Y_t Y_{t-2}] + E[Y_t \varepsilon_t] \\
\gamma_0 &= \phi_1 \gamma_1 + \phi_2 \gamma_2 + \sigma^2
\end{aligned}
\]

et donc la variance à long terme est une fonction des deux premières
autocovariances, des paramètres du modèle, et de la variance des
innovations. Cela peut être répété pour la première autocovariance,

\[
\begin{aligned}
E[Y_t Y_{t-1}] &= \phi_1 E[Y_{t-1} Y_{t-1}] + \phi_2 E[Y_{t-1} Y_{t-2}] + E[Y_{t-1} \varepsilon_t] \\
\gamma_1 &= \phi_1 \gamma_0 + \phi_2 \gamma_1,
\end{aligned}
\]

et pour la deuxième autocovariance,

\[
\begin{aligned}
E[Y_t Y_{t-2}] &= \phi_1 E[Y_{t-2} Y_{t-1}] + \phi_2 E[Y_{t-2} Y_{t-2}] + E[Y_{t-2} \varepsilon_t] \\
\gamma_2 &= \phi_1 \gamma_1 + \phi_2 \gamma_0.
\end{aligned}
\]

Ensemble, les trois dernières équations forment un système de trois
équations avec trois inconnues. La méthode de Yule-Walker repose
fortement sur la stationnarité de covariance et donc
\(E[Y_t Y_{t-j}] = E[Y_{t-h} Y_{t-h-j}]\) pour tout \(h\). Cette
propriété des processus stationnaires de covariance a été utilisée à
plusieurs reprises pour former les équations de Yule-Walker puisque
\(E[Y_t Y_t] = E[Y_{t-1} Y_{t-1}] = E[Y_{t-2} Y_{t-2}] = \gamma_0\) et
\(E[Y_t Y_{t-1}] = E[Y_{t-1} Y_{t-2}] = \gamma_1\). La méthode de
Yule-Walker sera démontrée pour plusieurs modèles, en commençant par un
simple MA(1) et en allant jusqu'à un ARMA(1,1).

\subsubsection{A.1.1 MA(1)}\label{a.1.1-ma1}

Les autocorrelations du MA(1) sont simples à dériver.

\[
Y_t = \theta_1 \varepsilon_{t-1} + \varepsilon_t
\]

Les équations de Yule-Walker sont

\[
\begin{aligned}
E[Y_t Y_t] &= E[\theta_1 \varepsilon_{t-1} Y_t] + E[\varepsilon_t Y_t] \\
E[Y_t Y_{t-1}] &= E[\theta_1 \varepsilon_{t-1} Y_{t-1}] + E[\varepsilon_t Y_{t-1}] \\
E[Y_t Y_{t-2}] &= E[\theta_1 \varepsilon_{t-1} Y_{t-2}] + E[\varepsilon_t Y_{t-2}]
\end{aligned}
\]

\[
\begin{aligned}
\gamma_0 &= \theta_1^2 \sigma^2 + \sigma^2 \\
\gamma_1 &= \theta_1 \sigma^2 \\
\gamma_2 &= 0
\end{aligned}
\]

De plus, \(\gamma_s\) et \(\rho_s\), pour \(s \geq 2\), sont nuls en
raison de la propriété de bruit blanc des résidus, et donc les
autocorrelations sont

\[
\rho_1 = \frac{\theta_1 \sigma^2}{\theta_1^2 \sigma^2 + \sigma^2} = \frac{\theta_1}{1 + \theta_1^2}
\]

\[
\rho_2 = 0.
\]

\subsubsection{A.1.2 MA(Q)}\label{a.1.2-maq}

Les équations de Yule-Walker peuvent être construites et résolues pour
tout MA(Q), et la structure de l'autocovariance est simple à détecter en
construisant un sous-ensemble du système complet.

\[
\begin{aligned}
E[Y_t Y_t] &= E[\theta_1 \varepsilon_{t-1} Y_t] + E[\theta_2 \varepsilon_{t-2} Y_t] + E[\theta_3 \varepsilon_{t-3} Y_t] + \ldots + E[\theta_Q \varepsilon_{t-Q} Y_t] \\
\gamma_0 &= \theta_1^2 \sigma^2 + \theta_2^2 \sigma^2 + \theta_3^2 \sigma^2 + \ldots + \theta_Q^2 \sigma^2 + \sigma^2 \\
&= \sigma^2 (1 + \theta_1^2 + \theta_2^2 + \theta_3^2 + \ldots + \theta_Q^2)
\end{aligned}
\]

\[
\begin{aligned}
E[Y_t Y_{t-1}] &= E[\theta_1 \varepsilon_{t-1} Y_{t-1}] + E[\theta_2 \varepsilon_{t-2} Y_{t-1}] + E[\theta_3 \varepsilon_{t-3} Y_{t-1}] + \ldots + E[\theta_Q \varepsilon_{t-Q} Y_{t-1}] \\
\gamma_1 &= \theta_1 \sigma^2 + \theta_1 \theta_2 \sigma^2 + \theta_2 \theta_3 \sigma^2 + \ldots + \theta_{Q-1} \theta_Q \sigma^2 \\
&= \sigma^2 (\theta_1 + \theta_1 \theta_2 + \theta_2 \theta_3 + \ldots + \theta_{Q-1} \theta_Q)
\end{aligned}
\]

\[
\begin{aligned}
E[Y_t Y_{t-2}] &= E[\theta_1 \varepsilon_{t-1} Y_{t-2}] + E[\theta_2 \varepsilon_{t-2} Y_{t-2}] + E[\theta_3 \varepsilon_{t-3} Y_{t-2}] + \ldots + E[\theta_Q \varepsilon_{t-Q} Y_{t-2}] \\
\gamma_2 &= \theta_2 \sigma^2 + \theta_1 \theta_3 \sigma^2 + \theta_2 \theta_4 \sigma^2 + \ldots + \theta_{Q-2} \theta_Q \sigma^2 \\
&= \sigma^2 (\theta_2 + \theta_1 \theta_3 + \theta_2 \theta_4 + \ldots + \theta_{Q-2} \theta_Q)
\end{aligned}
\]

Le schéma qui émerge montre que

\[
\gamma_s = \theta_s \sigma^2 + \sum_{i=1}^{Q-s} \theta_i \theta_{i+s} \sigma^2 = \sigma^2 \left( \theta_s + \sum_{i=1}^{Q-s} \theta_i \theta_{i+s} \right)
\]

et donc, \(\gamma_s\) est une somme de \(Q-s + 1\) termes. Les
autocorrélations sont

\[
\rho_1 = \frac{\theta_1 + \sum_{i=1}^{Q-1} \theta_i \theta_{i+1}}{1 + \sum_{i=1}^{Q} \theta_i^2}
\]

\[
\rho_2 = \frac{\theta_2 + \sum_{i=1}^{Q-2} \theta_i \theta_{i+2}}{1 + \sum_{i=1}^{Q} \theta_i^2}
\]

\[
\vdots
\]

\[
\rho_Q = \frac{\theta_Q}{1 + \sum_{i=1}^{Q} \theta_i^2}
\]

\[
\rho_{Q+s} = 0, \quad s \geq 0
\]

\subsubsection{A.1.3 AR(1)}\label{a.1.3-ar1}

La méthode de Yule-Walker nécessite \(\max(P, Q) + 1\) équations pour
calculer l'autocovariance d'un processus ARMA(P, Q), et dans le cas d'un
AR(1), deux équations sont requises (la troisième est incluse pour
établir ce point).

\[
Y_t = \phi_1 Y_{t-1} + \varepsilon_t
\]

\[
\begin{aligned}
E[Y_t Y_t] &= E[\phi_1 Y_{t-1} Y_t] + E[\varepsilon_t Y_t] \\
E[Y_t Y_{t-1}] &= E[\phi_1 Y_{t-1} Y_{t-1}] + E[\varepsilon_t Y_{t-1}] \\
E[Y_t Y_{t-2}] &= E[\phi_1 Y_{t-1} Y_{t-2}] + E[\varepsilon_t Y_{t-2}]
\end{aligned}
\]

Ces équations peuvent être réécrites en termes d'autocovariances, de
paramètres du modèle et de \(\sigma^2\) en prenant l'espérance et en
notant que \(E[\varepsilon_t Y_t] = \sigma^2\) puisque
\(Y_t = \varepsilon_t + \phi_1 \varepsilon_{t-1} + \phi_1^2 \varepsilon_{t-2} + \ldots\)
et \(E[\varepsilon_t Y_{t-j}] = 0\) pour \(j > 0\) car
\(\{\varepsilon_t\}\) est un processus de bruit blanc.

\[
\begin{aligned}
\gamma_0 &= \phi_1 \gamma_1 + \sigma^2 \\
\gamma_1 &= \phi_1 \gamma_0 \\
\gamma_2 &= \phi_1 \gamma_1
\end{aligned}
\]

La troisième équation est redondante car \(\gamma_2\) est entièrement
déterminé par \(\gamma_1\) et \(\phi_1\), et les autocovariances d'ordre
supérieur sont similairement redondantes puisque
\(\gamma_s = \phi_1 \gamma_{s-1}\) pour tout \(s\). Les deux premières
équations peuvent être résolues pour \(\gamma_0\) et \(\gamma_1\),

\[
\begin{aligned}
\gamma_0 &= \phi_1 \gamma_1 + \sigma^2 \\
\gamma_1 &= \phi_1 \gamma_0 \\
\Rightarrow \gamma_0 &= \phi_1^2 \gamma_0 + \sigma^2 \\
\Rightarrow \gamma_0 - \phi_1^2 \gamma_0 &= \sigma^2 \\
\Rightarrow \gamma_0 (1 - \phi_1^2) &= \sigma^2 \\
\Rightarrow \gamma_0 &= \frac{\sigma^2}{1 - \phi_1^2}
\end{aligned}
\]

et

\[
\begin{aligned}
\gamma_1 &= \phi_1 \gamma_0 \\
\Rightarrow \gamma_1 &= \phi_1 \frac{\sigma^2}{1 - \phi_1^2}.
\end{aligned}
\]

Les autocovariances restantes peuvent être calculées en utilisant la
récurrence \(\gamma_s = \phi_1 \gamma_{s-1}\), et donc

\[
\gamma_s = \phi_1^s \frac{\sigma^2}{1 - \phi_1^2}.
\]

Enfin, les autocorrélations peuvent être calculées comme des rapports
d'autocovariances,

\[
\rho_1 = \frac{\gamma_1}{\gamma_0} = \phi_1
\]

\[
\rho_s = \frac{\gamma_s}{\gamma_0} = \phi_1^s.
\]

\subsubsection{A.1.4 AR(2)}\label{a.1.4-ar2}

Les autocorrélations dans un AR(2)

\[
Y_t = \phi_1 Y_{t-1} + \phi_2 Y_{t-2} + \varepsilon_t
\]

peuvent être calculées de manière similaire en utilisant le système
d'équations de Yule-Walker avec \(\max(P, Q) + 1\) équations,

\[
\begin{aligned}
E[Y_t Y_t] &= \phi_1 E[Y_{t-1} Y_t] + \phi_2 E[Y_{t-2} Y_t] + E[\varepsilon_t Y_t] \\
E[Y_t Y_{t-1}] &= \phi_1 E[Y_{t-1} Y_{t-1}] + \phi_2 E[Y_{t-2} Y_{t-1}] + E[\varepsilon_t Y_{t-1}] \\
E[Y_t Y_{t-2}] &= \phi_1 E[Y_{t-1} Y_{t-2}] + \phi_2 E[Y_{t-2} Y_{t-2}] + E[\varepsilon_t Y_{t-2}]
\end{aligned}
\]

puis en remplaçant les espérances par leurs contreparties
populationnelles, \(\gamma_0\), \(\gamma_1\), \(\gamma_2\) et
\(\sigma^2\).

\[
\begin{aligned}
\gamma_0 &= \phi_1 \gamma_1 + \phi_2 \gamma_2 + \sigma^2 \\
\gamma_1 &= \phi_1 \gamma_0 + \phi_2 \gamma_1 \\
\gamma_2 &= \phi_1 \gamma_1 + \phi_2 \gamma_0
\end{aligned}
\]

De plus, il doit être vrai que
\(\gamma_s = \phi_1 \gamma_{s-1} + \phi_2 \gamma_{s-2}\) pour
\(s \geq 2\). Pour résoudre ce système d'équations, divisons les
équations d'autocovariance par \(\gamma_0\), la variance à long terme.
En omettant la première équation, le système se réduit à deux équations
avec deux inconnues,

\[
\begin{aligned}
\rho_1 &= \phi_1 \rho_0 + \phi_2 \rho_1 \\
\rho_2 &= \phi_1 \rho_1 + \phi_2 \rho_0
\end{aligned}
\]

puisque \(\rho_0 = \gamma_0 / \gamma_0 = 1\).

\[
\begin{aligned}
\rho_1 &= \phi_1 + \phi_2 \rho_1 \\
\rho_1 - \phi_2 \rho_1 &= \phi_1 \\
\rho_1 (1 - \phi_2) &= \phi_1 \\
\rho_1 &= \frac{\phi_1}{1 - \phi_2}
\end{aligned}
\]

et

\[
\begin{aligned}
\rho_2 &= \phi_1 \rho_1 + \phi_2 \\
&= \phi_1 \left( \frac{\phi_1}{1 - \phi_2} \right) + \phi_2 \\
&= \frac{\phi_1^2 + \phi_2 (1 - \phi_2)}{1 - \phi_2} \\
&= \frac{\phi_1^2 + \phi_2 - \phi_2^2}{1 - \phi_2}
\end{aligned}
\]

Puisque \(\rho_s = \phi_1 \rho_{s-1} + \phi_2 \rho_{s-2}\), ces deux
premières autocorrélations sont suffisantes pour calculer les autres
autocorrélations,

\[
\rho_3 = \phi_1 \rho_2 + \phi_2 \rho_1 = \frac{\phi_1 (\phi_1^2 + \phi_2 - \phi_2^2) + \phi_2 \phi_1}{1 - \phi_2}
\]

et la variance à long terme de \(Y_t\),

\[
\begin{aligned}
\gamma_0 &= \phi_1 \gamma_1 + \phi_2 \gamma_2 + \sigma^2 \\
\gamma_0 - \phi_1 \gamma_1 - \phi_2 \gamma_2 &= \sigma^2 \\
\gamma_0 (1 - \phi_1 \rho_1 - \phi_2 \rho_2) &= \sigma^2 \\
\gamma_0 &= \frac{\sigma^2}{1 - \phi_1 \rho_1 - \phi_2 \rho_2}
\end{aligned}
\]

La solution finale est calculée en substituant \(\rho_1\) et \(\rho_2\),

\[
\begin{aligned}
\gamma_0 &= \frac{\sigma^2}{1-\phi_1\frac{\phi_1}{1-\phi_2}-\phi_2\frac{\phi_1^2 + \phi_2 - \phi_2^2}{1 - \phi_2}}\\
         &= \frac{1 - \phi_2}{1 + \phi_2} \cdot \frac{\sigma^2}{(\phi_1 + \phi_2 - 1)(\phi_2 - \phi_1 - 1)}
\end{aligned}
\]

\subsubsection{A.1.5 AR(3)}\label{a.1.5-ar3}

Commencez par construire les équations de Yule-Walker, \[
\begin{aligned}
\mathbb{E}[Y_t Y_t] &= \phi_1 \mathbb{E}[Y_{t-1} Y_t] + \phi_2 \mathbb{E}[Y_{t-2} Y_t] + \phi_3 \mathbb{E}[Y_{t-3} Y_t] + \mathbb{E}[\varepsilon_t Y_t] \\
\mathbb{E}[Y_t Y_{t-1}] &= \phi_1 \mathbb{E}[Y_{t-1} Y_{t-1}] + \phi_2 \mathbb{E}[Y_{t-2} Y_{t-1}] + \phi_3 \mathbb{E}[Y_{t-3} Y_{t-1}] + \mathbb{E}[\varepsilon_t Y_{t-1}] \\
\mathbb{E}[Y_t Y_{t-2}] &= \phi_1 \mathbb{E}[Y_{t-1} Y_{t-2}] + \phi_2 \mathbb{E}[Y_{t-2} Y_{t-2}] + \phi_3 \mathbb{E}[Y_{t-3} Y_{t-2}] + \mathbb{E}[\varepsilon_t Y_{t-2}] \\
\mathbb{E}[Y_t Y_{t-3}] &= \phi_1 \mathbb{E}[Y_{t-1} Y_{t-3}] + \phi_2 \mathbb{E}[Y_{t-2} Y_{t-3}] + \phi_3 \mathbb{E}[Y_{t-3} Y_{t-3}] + \mathbb{E}[\varepsilon_t Y_{t-4}].
\end{aligned}
\]

En remplaçant les espérances par leurs valeurs populationnelles,
\(\gamma_0, \gamma_1, \dots\) et \(\sigma^2\), les équations de
Yule-Walker peuvent être réécrites comme suit : \[
\begin{aligned}
\gamma_0 &= \phi_1 \gamma_1 + \phi_2 \gamma_2 + \phi_3 \gamma_3 + \sigma^2 \\
\gamma_1 &= \phi_1 \gamma_0 + \phi_2 \gamma_1 + \phi_3 \gamma_2 \\
\gamma_2 &= \phi_1 \gamma_1 + \phi_2 \gamma_0 + \phi_3 \gamma_1 \\
\gamma_3 &= \phi_1 \gamma_2 + \phi_2 \gamma_1 + \phi_3 \gamma_0
\end{aligned}
\] et la relation récursive
\(\gamma_s = \phi_1 \gamma_{s-1} + \phi_2 \gamma_{s-2} + \phi_3 \gamma_{s-3}\)
peut être observée pour \(s \geq 3\).

En omettant la première condition et en divisant par \(\gamma_0\), \[
\begin{aligned}
\rho_1 &= \phi_1 \rho_0 + \phi_2 \rho_1 + \phi_3 \rho_2 \\
\rho_2 &= \phi_1 \rho_1 + \phi_2 \rho_0 + \phi_3 \rho_1 \\
\rho_3 &= \phi_1 \rho_2 + \phi_2 \rho_1 + \phi_3 \rho_0.
\end{aligned}
\] Cela laisse trois équations à trois inconnues puisque
\(\rho_0 = \gamma_0 / \gamma_0 = 1\).

\[
\begin{aligned}
\rho_1 &= \phi_1 + \phi_2 \rho_1 + \phi_3 \rho_2 \\
\rho_2 &= \phi_1 \rho_1 + \phi_2 + \phi_3 \rho_1 \\
\rho_3 &= \phi_1 \rho_2 + \phi_2 \rho_1 + \phi_3
\end{aligned}
\]

Après quelques calculs fastidieux, la solution à ce système est : \[
\begin{aligned}
\rho_1 &= \frac{\phi_1 + \phi_2 \phi_3}{1 - \phi_2 - \phi_1 \phi_3 - \phi_3^2} \\
\rho_2 &= \frac{\phi_2 + \phi_1^2 + \phi_3 \phi_1 - \phi_2^2}{1 - \phi_2 - \phi_1 \phi_3 - \phi_3^2} \\
\rho_3 &= \frac{\phi_3 + \phi_1^3 + \phi_1^2 \phi_3 + \phi_1 \phi_2^2 + 2 \phi_1 \phi_2 + \phi_2^2 \phi_3 - \phi_2 \phi_3 - \phi_1 \phi_3^2 - \phi_3^3}{1 - \phi_2 - \phi_1 \phi_3 - \phi_3^2}.
\end{aligned}
\]

Enfin, la variance inconditionnelle peut être calculée en utilisant les
trois premières autocorrélations, \[
\begin{aligned}
\gamma_0 &= \phi_1 \gamma_1 + \phi_2 \gamma_2 + \phi_3 \gamma_3 + \sigma^2 \\
\gamma_0 - \phi_1 \gamma_1 - \phi_2 \gamma_2 - \phi_3 \gamma_3 &= \sigma^2 \\
\gamma_0 (1 - \phi_1 \rho_1 - \phi_2 \rho_2 - \phi_3 \rho_3) &= \sigma^2 \\
\gamma_0 &= \frac{\sigma^2}{1 - \phi_1 \rho_1 - \phi_2 \rho_2 - \phi_3 \rho_3} \\
\gamma_0 &= \frac{\sigma^2 (1 - \phi_2 - \phi_1 \phi_3 - \phi_3^2)}{(1 - \phi_2 - \phi_3 - \phi_1)(1 + \phi_2 + \phi_3 \phi_1 - \phi_3^2)(1 + \phi_3 + \phi_1 - \phi_2)}.
\end{aligned}
\]

\subsubsection{A.1.6 ARMA(1,1)}\label{a.1.6-arma11}

Le calcul des autocovariances et autocorrélations d'un processus ARMA
est plus complexe que pour un processus AR ou MA pur. Un ARMA(1,1) est
spécifié comme suit : \[
Y_t = \phi_1 Y_{t-1} + \theta_1 \varepsilon_{t-1} + \varepsilon_t
\] Puisque ( P = Q = 1 ), le système de Yule-Walker nécessite deux
équations, en notant que l'autocovariance d'ordre 3 ou plus est une
fonction triviale des deux premières autocovariances.

\[
\begin{aligned}
\mathbb{E}[Y_t Y_t] &= \mathbb{E}[\phi_1 Y_{t-1} Y_t] + \mathbb{E}[\theta_1 \varepsilon_{t-1} Y_t] + \mathbb{E}[\varepsilon_t Y_t] \\
\mathbb{E}[Y_t Y_{t-1}] &= \mathbb{E}[\phi_1 Y_{t-1} Y_{t-1}] + \mathbb{E}[\theta_1 \varepsilon_{t-1} Y_{t-1}] + \mathbb{E}[\varepsilon_t Y_{t-1}]
\end{aligned}
\]

La présence du terme ( \mathbb{E}{[}\theta\emph{1 \varepsilon}\{t-1\}
Y\_t{]} ) dans la première équation complique la résolution du système,
car ( \varepsilon\emph{\{t-1\} ) apparaît directement dans ( Y\_t ) via
( \theta\emph{1 \varepsilon}\{t-1\} ) et indirectement via ( \phi\emph{1
Y}\{t-1\} ). Les relations non nulles peuvent être déterminées en
substituant récursivement ( Y\_t ) jusqu'à ce qu'il ne consiste qu'en (
\varepsilon\emph{t ), ( \varepsilon}\{t-1\} ) et ( Y}\{t-2\} ) (puisque
( Y\_\{t-2\} ) est non corrélé avec ( \varepsilon\_\{t-1\} ) par
l'hypothèse de bruit blanc).

\[
\begin{aligned}
Y_t &= \phi_1 Y_{t-1} + \theta_1 \varepsilon_{t-1} + \varepsilon_t \\
&= \phi_1 (\phi_1 Y_{t-2} + \theta_1 \varepsilon_{t-2} + \varepsilon_{t-1}) + \theta_1 \varepsilon_{t-1} + \varepsilon_t \\
&= \phi_1^2 Y_{t-2} + \phi_1 \theta_1 \varepsilon_{t-2} + \phi_1 \varepsilon_{t-1} + \theta_1 \varepsilon_{t-1} + \varepsilon_t \\
&= \phi_1^2 Y_{t-2} + \phi_1 \theta_1 \varepsilon_{t-2} + (\phi_1 + \theta_1) \varepsilon_{t-1} + \varepsilon_t
\end{aligned}
\]

Ainsi, ( \mathbb{E}{[}\theta\emph{1 \varepsilon}\{t-1\} Y\_t{]} =
\theta\_1 (\phi\_1 + \theta\_1) \sigma\^{}2 ), et les équations de
Yule-Walker peuvent être exprimées en utilisant les moments
populationnels et les paramètres du modèle.

\[
\begin{aligned}
\gamma_0 &= \phi_1 \gamma_1 + \theta_1 (\phi_1 + \theta_1) \sigma^2 + \sigma^2 \\
\gamma_1 &= \phi_1 \gamma_0 + \theta_1 \sigma^2
\end{aligned}
\]

Ces deux équations à deux inconnues peuvent être résolues :

\[
\begin{aligned}
\gamma_0 &= \phi_1 \gamma_1 + \theta_1 (\phi_1 + \theta_1) \sigma^2 + \sigma^2 \\
&= \phi_1 (\phi_1 \gamma_0 + \theta_1 \sigma^2) + \theta_1 (\phi_1 + \theta_1) \sigma^2 + \sigma^2 \\
&= \phi_1^2 \gamma_0 + \phi_1 \theta_1 \sigma^2 + \theta_1 (\phi_1 + \theta_1) \sigma^2 + \sigma^2 \\
\gamma_0 - \phi_1^2 \gamma_0 &= \sigma^2 (\phi_1 \theta_1 + \phi_1 \theta_1 + \theta_1^2 + 1) \\
\gamma_0 &= \frac{\sigma^2 (1 + \theta_1^2 + 2 \phi_1 \theta_1)}{1 - \phi_1^2}
\end{aligned}
\]

De même, pour ( \gamma\_1 ) :

\[
\begin{aligned}
\gamma_1 &= \phi_1 \gamma_0 + \theta_1 \sigma^2 \\
&= \phi_1 \left( \frac{\sigma^2 (1 + \theta_1^2 + 2 \phi_1 \theta_1)}{1 - \phi_1^2} \right) + \theta_1 \sigma^2 \\
&= \frac{\sigma^2 (\phi_1 + \phi_1 \theta_1^2 + 2 \phi_1^2 \theta_1)}{1 - \phi_1^2} + \theta_1 \sigma^2 \\
&= \frac{\sigma^2 (\phi_1 + \phi_1 \theta_1^2 + 2 \phi_1^2 \theta_1) + \theta_1 \sigma^2 (1 - \phi_1^2)}{1 - \phi_1^2} \\
&= \frac{\sigma^2 (\phi_1 + \phi_1 \theta_1^2 + 2 \phi_1^2 \theta_1 + \theta_1 - \phi_1^2 \theta_1)}{1 - \phi_1^2} \\
&= \frac{\sigma^2 (\phi_1 + \theta_1 + \phi_1 \theta_1^2 + \phi_1^2 \theta_1)}{1 - \phi_1^2} \\
&= \frac{\sigma^2 (\phi_1 + \theta_1)(1 + \phi_1 \theta_1)}{1 - \phi_1^2} \\
\end{aligned}
\] Ainsi, la première autocorrélation est :

\[
\rho_1 = \frac{\gamma_1}{\gamma_0} = \frac{\frac{\sigma^2 (\phi_1 + \theta_1)(1 + \phi_1 \theta_1)}{1 - \phi_1^2}}{\frac{\sigma^2 (1 + \theta_1^2 + 2 \phi_1 \theta_1)}{1 - \phi_1^2}} = \frac{(\phi_1 + \theta_1)(1 + \phi_1 \theta_1)}{1 + \theta_1^2 + 2 \phi_1 \theta_1}
\]

En revenant à l'équation suivante de Yule-Walker :

\[
\mathbb{E}[Y_t Y_{t-2}] = \mathbb{E}[\phi_1 Y_{t-1} Y_{t-2}] + \mathbb{E}[\theta_1 \varepsilon_{t-1} Y_{t-2}] + \mathbb{E}[\varepsilon_t Y_{t-2}]
\]

Ainsi, ( \gamma\_2 = \phi\_1 \gamma\_1 ), et en divisant les deux côtés
par ( \gamma\_0 ), on obtient ( \rho\_2 = \phi\_1 \rho\_1 ). Les
autocovariances et autocorrélations d'ordre supérieur suivent
respectivement ( \gamma\_s = \phi\emph{1 \gamma}\{s-1\} ) et ( \rho\_s =
\phi\emph{1 \rho}\{s-1\} ), et donc ( \rho\_s = \phi\_1\^{}\{s-1\}
\rho\_1 ) pour ( s \geq 2 ).

\subsection{A.2 Substitution
rétroactive}\label{a.2-substitution-ruxe9troactive}

La substitution rétroactive est une méthode directe mais fastidieuse
pour dériver la fonction d'autocorrélation (ACF) et la variance à long
terme.

\subsubsection{A.2.1 AR(1)}\label{a.2.1-ar1}

Le processus AR(1), \[
Y_t = \phi_1 Y_{t-1} + \varepsilon_t,
\] est stationnaire si (\textbar{}\phi\_1\textbar{} \textless{} 1) et si
(\{\varepsilon\_t\}) est un bruit blanc. Pour calculer les
autocovariances et autocorrélations en utilisant la substitution
rétroactive, on transforme (Y\_t = \phi\emph{1 Y}\{t-1\} +
\varepsilon\_t) en un processus MA pur par substitution récursive :

\[
\begin{aligned}
Y_t &= \phi_1 Y_{t-1} + \varepsilon_t  \\
&= \phi_1 (\phi_1 Y_{t-2} + \varepsilon_{t-1}) + \varepsilon_t \\
&= \phi_1^2 Y_{t-2} + \phi_1 \varepsilon_{t-1} + \varepsilon_t \\
&= \phi_1^2 (\phi_1 Y_{t-3} + \varepsilon_{t-2}) + \phi_1 \varepsilon_{t-1} + \varepsilon_t \\
&= \phi_1^3 Y_{t-3} + \phi_1^2 \varepsilon_{t-2} + \phi_1 \varepsilon_{t-1} + \varepsilon_t \\
&= \varepsilon_t + \phi_1 \varepsilon_{t-1} + \phi_1^2 \varepsilon_{t-2} + \phi_1^3 \varepsilon_{t-3} + \dots \\
Y_t &= \sum_{i=0}^{\infty} \phi_1^i \varepsilon_{t-i}.
\end{aligned}
\]

La variance est l'espérance du carré : \[
\begin{aligned}
\gamma_0 &= \text{Var}[Y_t] = \mathbb{E}[Y_t^2]  \\
&= \mathbb{E}\left[\left(\sum_{i=0}^{\infty} \phi_1^i \varepsilon_{t-i}\right)^2\right] \\
&= \mathbb{E}\left[(\varepsilon_t + \phi_1 \varepsilon_{t-1} + \phi_1^2 \varepsilon_{t-2} + \phi_1^3 \varepsilon_{t-3} + \dots)^2\right] \\
&= \mathbb{E}\left[\sum_{i=0}^{\infty} \phi_1^{2i} \varepsilon_{t-i}^2 + \sum_{i=0}^{\infty} \sum_{j=0, j \neq i}^{\infty} \phi_1^i \phi_1^j \varepsilon_{t-i} \varepsilon_{t-j}\right] \\
&= \sum_{i=0}^{\infty} \phi_1^{2i} \mathbb{E}[\varepsilon_{t-i}^2] + \sum_{i=0}^{\infty} \sum_{j=0, j \neq i}^{\infty} \phi_1^i \phi_1^j \mathbb{E}[\varepsilon_{t-i} \varepsilon_{t-j}] \\
&= \sum_{i=0}^{\infty} \phi_1^{2i} \sigma^2 + \sum_{i=0}^{\infty} \sum_{j=0, j \neq i}^{\infty} \phi_1^i \phi_1^j \cdot 0 \\
&= \sum_{i=0}^{\infty} \phi_1^{2i} \sigma^2 \\
&= \frac{\sigma^2}{1 - \phi_1^2}.
\end{aligned}
\]

L'étape difficile dans la dérivation est la séparation des
(\varepsilon\emph{\{t-i\}) en ceux qui sont appariés à leur propre
retard ((\varepsilon}\{t-i\}\^{}2)) et ceux qui ne le sont pas
((\varepsilon\emph{\{t-i\} \varepsilon}\{t-j\}, i \neq j)). Le reste de
la dérivation découle de l'hypothèse que (\{\varepsilon\emph{t\}) est un
bruit blanc, donc (\mathbb{E}{[}\varepsilon\_\{t-i\}\^{}2{]} =
\sigma\^{}2) et (\mathbb{E}{[}\varepsilon\emph{\{t-i\}
\varepsilon}\{t-j\}{]} = 0) pour (i \neq j). Enfin, l'identité (\lim}\{n
\to \infty\} \sum\_\{i=0\}\^{}n \phi\_1\^{}\{2i\} =
\frac{1}{1 - \phi_1^2}) pour (\textbar{}\phi\_1\textbar{} \textless{} 1)
a été utilisée pour simplifier l'expression.

La première autocovariance peut être calculée en utilisant les mêmes
étapes sur la représentation MA((\infty)) : \[
\begin{aligned}
\gamma_1 &= \mathbb{E}[Y_t Y_{t-1}]  \\
&= \mathbb{E}\left[\sum_{i=0}^{\infty} \phi_1^i \varepsilon_{t-i} \sum_{j=1}^{\infty} \phi_1^{j-1} \varepsilon_{t-j}\right] \\
&= \mathbb{E}\left[(\varepsilon_t + \phi_1 \varepsilon_{t-1} + \phi_1^2 \varepsilon_{t-2} + \dots)(\varepsilon_{t-1} + \phi_1 \varepsilon_{t-2} + \phi_1^2 \varepsilon_{t-3} + \dots)\right] \\
&= \mathbb{E}\left[\phi_1 \sum_{i=0}^{\infty} \phi_1^{2i} \varepsilon_{t-1-i}^2 + \sum_{i=0}^{\infty} \sum_{j=1, j \neq i}^{\infty} \phi_1^i \phi_1^{j-1} \varepsilon_{t-i} \varepsilon_{t-j}\right] \\
&= \phi_1 \sum_{i=0}^{\infty} \phi_1^{2i} \mathbb{E}[\varepsilon_{t-1-i}^2] + \sum_{i=0}^{\infty} \sum_{j=1, j \neq i}^{\infty} \phi_1^i \phi_1^{j-1} \mathbb{E}[\varepsilon_{t-i} \varepsilon_{t-j}] \\
&= \phi_1 \sum_{i=0}^{\infty} \phi_1^{2i} \sigma^2 + \sum_{i=0}^{\infty} \sum_{j=1, j \neq i}^{\infty} \phi_1^i \phi_1^{j-1} \cdot 0 \\
&= \phi_1 \sum_{i=0}^{\infty} \phi_1^{2i} \sigma^2 \\
&= \phi_1 \gamma_0.
\end{aligned}
\]

De même, l'autocovariance d'ordre (s) peut être déterminée : \[
\begin{aligned}
\gamma_s &= \mathbb{E}[Y_t Y_{t-s}]  \\
&= \mathbb{E}\left[\sum_{i=0}^{\infty} \phi_1^i \varepsilon_{t-i} \sum_{j=s}^{\infty} \phi_1^{j-s} \varepsilon_{t-j}\right] \\
&= \mathbb{E}\left[\phi_1^s \sum_{i=0}^{\infty} \phi_1^{2i} \varepsilon_{t-s-i}^2 + \sum_{i=0}^{\infty} \sum_{j=s, j \neq i}^{\infty} \phi_1^i \phi_1^{j-s} \varepsilon_{t-i} \varepsilon_{t-j}\right] \\
&= \phi_1^s \sum_{i=0}^{\infty} \phi_1^{2i} \mathbb{E}[\varepsilon_{t-s-i}^2] + \sum_{i=0}^{\infty} \sum_{j=s, j \neq i}^{\infty} \phi_1^i \phi_1^{j-s} \mathbb{E}[\varepsilon_{t-i} \varepsilon_{t-j}] \\
&= \phi_1^s \sum_{i=0}^{\infty} \phi_1^{2i} \sigma^2 + \sum_{i=0}^{\infty} \sum_{j=s, j \neq i}^{\infty} \phi_1^i \phi_1^{j-s} \cdot 0 \\
&= \phi_1^s \sum_{i=0}^{\infty} \phi_1^{2i} \sigma^2 \\
&= \phi_1^s \gamma_0.
\end{aligned}
\]

Enfin, les autocorrélations peuvent être calculées à partir des ratios
des autocovariances : \[
\rho_1 = \frac{\gamma_1}{\gamma_0} = \phi_1 \quad \text{et} \quad \rho_s = \frac{\gamma_s}{\gamma_0} = \phi_1^s.
\]

\subsubsection{A.2.2 MA(1)}\label{a.2.2-ma1}

Le modèle MA(1) est le modèle de série temporelle non dégénéré le plus
simple considéré dans ce cours : {[} Y\_t = \theta\emph{1
\varepsilon}\{t-1\} + \varepsilon\_t. {]}

La dérivation de sa fonction d'autocorrélation est triviale, car aucune
substitution rétroactive n'est nécessaire. La variance est donnée par :
{[}

\begin{aligned}
\gamma_0 &= \text{Var}[Y_t] = \mathbb{E}[Y_t^2]  \\
&= \mathbb{E}[(\theta_1 \varepsilon_{t-1} + \varepsilon_t)^2] \\
&= \mathbb{E}[\theta_1^2 \varepsilon_{t-1}^2 + 2 \theta_1 \varepsilon_t \varepsilon_{t-1} + \varepsilon_t^2] \\
&= \mathbb{E}[\theta_1^2 \varepsilon_{t-1}^2] + \mathbb{E}[2 \theta_1 \varepsilon_t \varepsilon_{t-1}] + \mathbb{E}[\varepsilon_t^2] \\
&= \theta_1^2 \sigma^2 + 0 + \sigma^2 \\
&= \sigma^2 (1 + \theta_1^2).
\end{aligned}

{]}

La première autocovariance est : {[}

\begin{aligned}
\gamma_1 &= \mathbb{E}[Y_t Y_{t-1}]  \\
&= \mathbb{E}[(\theta_1 \varepsilon_{t-1} + \varepsilon_t)(\theta_1 \varepsilon_{t-2} + \varepsilon_{t-1})] \\
&= \mathbb{E}[\theta_1^2 \varepsilon_{t-1} \varepsilon_{t-2} + \theta_1 \varepsilon_{t-1}^2 + \theta_1 \varepsilon_t \varepsilon_{t-2} + \varepsilon_t \varepsilon_{t-1}] \\
&= \mathbb{E}[\theta_1^2 \varepsilon_{t-1} \varepsilon_{t-2}] + \mathbb{E}[\theta_1 \varepsilon_{t-1}^2] + \mathbb{E}[\theta_1 \varepsilon_t \varepsilon_{t-2}] + \mathbb{E}[\varepsilon_t \varepsilon_{t-1}] \\
&= 0 + \theta_1 \sigma^2 + 0 + 0 \\
&= \theta_1 \sigma^2.
\end{aligned}

{]}

La deuxième autocovariance (et les autocovariances d'ordre supérieur)
est : {[}

\begin{aligned}
\gamma_2 &= \mathbb{E}[Y_t Y_{t-2}]  \\
&= \mathbb{E}[(\theta_1 \varepsilon_{t-1} + \varepsilon_t)(\theta_1 \varepsilon_{t-3} + \varepsilon_{t-2})] \\
&= \mathbb{E}[\theta_1^2 \varepsilon_{t-1} \varepsilon_{t-3} + \theta_1 \varepsilon_{t-1} \varepsilon_{t-2} + \theta_1 \varepsilon_t \varepsilon_{t-3} + \varepsilon_t \varepsilon_{t-2}] \\
&= \mathbb{E}[\theta_1^2 \varepsilon_{t-1} \varepsilon_{t-3}] + \mathbb{E}[\theta_1 \varepsilon_{t-1} \varepsilon_{t-2}] + \mathbb{E}[\theta_1 \varepsilon_t \varepsilon_{t-3}] + \mathbb{E}[\varepsilon_t \varepsilon_{t-2}] \\
&= 0 + 0 + 0 + 0 \\
&= 0.
\end{aligned}

{]}

Les autocorrélations sont donc : {[} \rho\_1 =
\frac{\theta_1}{1 + \theta_1^2}, \quad \rho\_s = 0 \quad \text{pour}
\quad s \geq 2. {]}

\subsubsection{A.2.3 ARMA(1,1)}\label{a.2.3-arma11}

Un processus ARMA(1,1) est défini par : {[} Y\_t = \phi\emph{1 Y}\{t-1\}
+ \theta\emph{1 \varepsilon}\{t-1\} + \varepsilon\_t. {]} Il est
stationnaire si (\textbar{}\phi\_1\textbar{} \textless{} 1) et si
(\{\varepsilon\_t\}) est un bruit blanc. La dérivation de la variance et
des autocovariances est plus complexe que pour un processus AR(1). Il
est à noter que cette dérivation est plus longue et plus complexe que la
résolution des équations de Yule-Walker.

\subsubsection{\texorpdfstring{Représentation
MA((\infty))}{Représentation MA(())}}\label{repruxe9sentation-ma}

On commence par calculer la représentation MA((\infty)) en utilisant la
substitution rétroactive : {[}

\begin{aligned}
Y_t &= \phi_1 Y_{t-1} + \theta_1 \varepsilon_{t-1} + \varepsilon_t \\
&= \phi_1 (\phi_1 Y_{t-2} + \theta_1 \varepsilon_{t-2} + \varepsilon_{t-1}) + \theta_1 \varepsilon_{t-1} + \varepsilon_t \\
&= \phi_1^2 Y_{t-2} + \phi_1 \theta_1 \varepsilon_{t-2} + \phi_1 \varepsilon_{t-1} + \theta_1 \varepsilon_{t-1} + \varepsilon_t \\
&= \phi_1^2 (\phi_1 Y_{t-3} + \theta_1 \varepsilon_{t-3} + \varepsilon_{t-2}) + \phi_1 \theta_1 \varepsilon_{t-2} + (\phi_1 + \theta_1) \varepsilon_{t-1} + \varepsilon_t \\
&= \phi_1^3 Y_{t-3} + \phi_1^2 \theta_1 \varepsilon_{t-3} + \phi_1^2 \varepsilon_{t-2} + \phi_1 \theta_1 \varepsilon_{t-2} + (\phi_1 + \theta_1) \varepsilon_{t-1} + \varepsilon_t \\
&= \phi_1^3 (\phi_1 Y_{t-4} + \theta_1 \varepsilon_{t-4} + \varepsilon_{t-3}) + \phi_1^2 \theta_1 \varepsilon_{t-3} + \phi_1 (\phi_1 + \theta_1) \varepsilon_{t-2} + (\phi_1 + \theta_1) \varepsilon_{t-1} + \varepsilon_t \\
&= \phi_1^4 Y_{t-4} + \phi_1^3 \theta_1 \varepsilon_{t-4} + \phi_1^3 \varepsilon_{t-3} + \phi_1^2 \theta_1 \varepsilon_{t-3} + \phi_1 (\phi_1 + \theta_1) \varepsilon_{t-2} + (\phi_1 + \theta_1) \varepsilon_{t-1} + \varepsilon_t \\
&= \phi_1^4 Y_{t-4} + \phi_1^3 \theta_1 \varepsilon_{t-4} + \phi_1^2 (\phi_1 + \theta_1) \varepsilon_{t-3} + \phi_1 (\phi_1 + \theta_1) \varepsilon_{t-2} + (\phi_1 + \theta_1) \varepsilon_{t-1} + \varepsilon_t \\
&= \varepsilon_t + (\phi_1 + \theta_1) \varepsilon_{t-1} + \phi_1 (\phi_1 + \theta_1) \varepsilon_{t-2} + \phi_1^2 (\phi_1 + \theta_1) \varepsilon_{t-3} + \dots \\
&= \varepsilon_t + \sum_{i=0}^{\infty} \phi_1^i (\phi_1 + \theta_1) \varepsilon_{t-1-i}. 
\end{aligned}

{]}

\subsubsection{\texorpdfstring{Variance
((\gamma\_0))}{Variance ((\_0))}}\label{variance-_0}

La variance est donnée par : {[}

\begin{aligned}
\gamma_0 &= \text{Var}[Y_t] = \mathbb{E}[Y_t^2] \\
&= \mathbb{E}\left[\left(\varepsilon_t + \sum_{i=0}^{\infty} \phi_1^i (\phi_1 + \theta_1) \varepsilon_{t-1-i}\right)^2\right] \\
&= \mathbb{E}\left[\varepsilon_t^2 + 2 \varepsilon_t \sum_{i=0}^{\infty} \phi_1^i (\phi_1 + \theta_1) \varepsilon_{t-1-i} + \left(\sum_{i=0}^{\infty} \phi_1^i (\phi_1 + \theta_1) \varepsilon_{t-1-i}\right)^2\right] \\
&= \mathbb{E}[\varepsilon_t^2] + \mathbb{E}\left[2 \varepsilon_t \sum_{i=0}^{\infty} \phi_1^i (\phi_1 + \theta_1) \varepsilon_{t-1-i}\right] + \mathbb{E}\left[\left(\sum_{i=0}^{\infty} \phi_1^i (\phi_1 + \theta_1) \varepsilon_{t-1-i}\right)^2\right] \\
&= \sigma^2 + 0 + \mathbb{E}\left[\sum_{i=0}^{\infty} \phi_1^{2i} (\phi_1 + \theta_1)^2 \varepsilon_{t-1-i}^2 + \sum_{i=0}^{\infty} \sum_{j=0, j \neq i}^{\infty} \phi_1^i \phi_1^j (\phi_1 + \theta_1)^2 \varepsilon_{t-1-i} \varepsilon_{t-1-j}\right] \\
&= \sigma^2 + \sum_{i=0}^{\infty} \phi_1^{2i} (\phi_1 + \theta_1)^2 \sigma^2 + \sum_{i=0}^{\infty} \sum_{j=0, j \neq i}^{\infty} \phi_1^i \phi_1^j (\phi_1 + \theta_1)^2 \cdot 0 \\
&= \sigma^2 + (\phi_1 + \theta_1)^2 \sigma^2 \sum_{i=0}^{\infty} \phi_1^{2i} \\
&= \sigma^2 + (\phi_1 + \theta_1)^2 \sigma^2 \frac{1}{1 - \phi_1^2} \\
&= \sigma^2 \left(1 + \frac{(\phi_1 + \theta_1)^2}{1 - \phi_1^2}\right) \\
&= \sigma^2 \frac{1 - \phi_1^2 + (\phi_1 + \theta_1)^2}{1 - \phi_1^2} \\
&= \sigma^2 \frac{1 + \theta_1^2 + 2 \phi_1 \theta_1}{1 - \phi_1^2}.
\end{aligned}

{]}

\subsubsection{\texorpdfstring{Autocovariance d'ordre 1
((\gamma\_1))}{Autocovariance d'ordre 1 ((\_1))}}\label{autocovariance-dordre-1-_1}

L'autocovariance d'ordre 1 est donnée par : {[}

\begin{aligned}
\gamma_1 &= \mathbb{E}[Y_t Y_{t-1}] \\
&= \mathbb{E}\left[\left(\varepsilon_t + \sum_{i=0}^{\infty} \phi_1^i (\phi_1 + \theta_1) \varepsilon_{t-1-i}\right) \left(\varepsilon_{t-1} + \sum_{i=0}^{\infty} \phi_1^i (\phi_1 + \theta_1) \varepsilon_{t-2-i}\right)\right] \\
&= \mathbb{E}\left[\varepsilon_t \varepsilon_{t-1} + \sum_{i=0}^{\infty} \phi_1^i (\phi_1 + \theta_1) \varepsilon_t \varepsilon_{t-2-i} + \sum_{i=0}^{\infty} \phi_1^i (\phi_1 + \theta_1) \varepsilon_{t-1} \varepsilon_{t-1-i} \right. \\
&\quad \left. + \left(\sum_{i=0}^{\infty} \phi_1^i (\phi_1 + \theta_1) \varepsilon_{t-1-i}\right) \left(\sum_{i=0}^{\infty} \phi_1^i (\phi_1 + \theta_1) \varepsilon_{t-2-i}\right)\right] \\
&= 0 + 0 + (\phi_1 + \theta_1) \sigma^2 + \mathbb{E}\left[\sum_{i=0}^{\infty} \phi_1^{2i+1} (\phi_1 + \theta_1)^2 \varepsilon_{t-2-i}^2 + \sum_{i=0}^{\infty} \sum_{j=0, j \neq i+1}^{\infty} \phi_1^i \phi_1^j (\phi_1 + \theta_1)^2 \varepsilon_{t-1-i} \varepsilon_{t-2-j}\right] \\
&= (\phi_1 + \theta_1) \sigma^2 + \phi_1 (\phi_1 + \theta_1)^2 \sigma^2 \sum_{i=0}^{\infty} \phi_1^{2i} \\
&= (\phi_1 + \theta_1) \sigma^2 + \phi_1 (\phi_1 + \theta_1)^2 \sigma^2 \frac{1}{1 - \phi_1^2} \\
&= \sigma^2 \left((\phi_1 + \theta_1) + \frac{\phi_1 (\phi_1 + \theta_1)^2}{1 - \phi_1^2}\right) \\
&= \sigma^2 \frac{(\phi_1 + \theta_1)(1 - \phi_1^2) + \phi_1 (\phi_1 + \theta_1)^2}{1 - \phi_1^2} \\
&= \sigma^2 \frac{(\phi_1 + \theta_1)(1 + \phi_1 \theta_1)}{1 - \phi_1^2}.
\end{aligned}

{]}

\subsubsection{Autocorrélations}\label{autocorruxe9lations}

L'autocorrélation d'ordre 1 est : {[} \rho\_1 =
\frac{\gamma_1}{\gamma_0} =
\frac{(\phi_1 + \theta_1)(1 + \phi_1 \theta_1)}{1 + \theta_1^2 + 2 \phi_1 \theta_1}.
{]}

Les autocorrélations d'ordre supérieur suivent la récurrence : {[}
\rho\_s = \phi\emph{1 \rho}\{s-1\}, \quad s \geq 2. {]}




\end{document}
